% ============================================================================
\chapter{Stockage, cache et fonctionnement hors ligne}
\label{ch:js-storage}
% ============================================================================

\begin{objectif}
\begin{itemize}
  \item Choisir le bon stockage selon quatre crit\`eres~: taille, dur\'ee, origine, secret.
  \item Utiliser \texttt{localStorage}/\texttt{sessionStorage} sans se faire pi\'eger par le JSON.
  \item Lire/\'ecrire des cookies correctement (expire, SameSite, Secure, HttpOnly).
  \item Manipuler IndexedDB proprement, y compris avec des index et des curseurs.
  \item Poser un service worker pour un vrai mode hors ligne et une PWA installable.
\end{itemize}
\end{objectif}

\section{Quoi stocker o\`u}

\begin{tabz}{llllX}
  \headrow \thd{M\'ecanisme} \& \thd{Capacit\'e typique} \& \thd{Dur\'ee} \& \thd{Serveur ?} \& \thd{Usage correct} \\
  \midrule
  cookie \& 4~Ko par domaine \& jusqu'\`a \verb|Expires| \& oui, \`a chaque requ\^ete \& session, pr\'ef\'erence lue c\^ot\'e serveur \\
  \verb|localStorage| \& 5--10~Mo \& infini, par origine \& non \& pr\'ef\'erences, brouillon de formulaire, \'etat d'UI \\
  \verb|sessionStorage| \& idem \& onglet courant \& non \& \'etat d'un wizard, jeton court \\
  \verb|IndexedDB| \& centaines de Mo (jusqu'au quota disque) \& infini \& non \& grandes quantit\'es structur\'ees, recherche, hors ligne \\
  OPFS (files) \& quota disque \& infini \& non \& fichiers, bases SQLite WASM, gros exports \\
  Cache API \& quota \& infini \& non \& r\'eponses HTTP mises en cache par le service worker \\
  \verb|caches|+SW \& --- \& --- \& --- \& application enti\`ere hors ligne \\
\end{tabz}

\begin{piege}
\textbf{Aucun stockage c\^ot\'e navigateur n'est digne de confiance.} L'utilisateur peut lire,
modifier, supprimer n'importe quoi (console, extension, disque). Ne stocke donc jamais l\`a~:
un r\^ole \og ~admin~ \fg{}, un prix, un solde, une validation de paiement, un jeton de session sensible
sans expiration courte. Le stockage local est un \emph{cache}, pas une source de v\'erit\'e.
\end{piege}

\section{Web Storage --- proprement, une fois pour toutes}

\begin{codejs}
// localStorage ne stocke QUE des chaines : un objet passe par JSON, et ca peut echouer.
class Stockage {
  constructor(prefixe = "app:", storage = localStorage) { Object.assign(this, { prefixe, storage }); }
  #cle(k) { return this.prefixe + k; }
  lire(k, defaut) {
    try { const brut = this.storage.getItem(this.#cle(k)); return brut === null ? defaut : JSON.parse(brut); }
    catch (e) { console.warn(`storage.lire(${k})`, e); this.effacer(k); return defaut; }
  }
  ecrire(k, valeur) {
    try { this.storage.setItem(this.#cle(k), JSON.stringify(valeur)); return true; }
    catch (e) {
      if (e.name === "QuotaExceededError") { this.nettoyer(); return this.ecrire(k, valeur); }
      console.warn(`storage.ecrire(${k})`, e); return false;
    }
  }
  effacer(k) { this.storage.removeItem(this.#cle(k)); }
  nettoyer() {  // enleve les cles trop vieilles ( convention : cle + "__t" = timestamp)
    const ages = Object.keys(this.storage)
      .filter((k) => k.startsWith(this.prefixe))
      .map((k) => ({ k, t: Number(this.storage.getItem(`${k}__t`) ?? 0) }))
      .sort((a, b) => a.t - b.t);
    ages.slice(0, Math.ceil(ages.length / 2)).forEach(({ k }) => this.effacer(k.slice(this.prefixe.length)));
  }
  get taille() { let n = 0; for (const k of Object.keys(this.storage)) if (k.startsWith(this.prefixe)) n += this.storage.getItem(k).length; return n; }
}
export const store = new Stockage("carnet:");
store.ecrire("articles", [{ id: 1 }]);
store.lire("articles", []);

// Synchro entre onglets : l'evenement "storage" se declenche sur les AUTRES onglets
window.addEventListener("storage", ({ key, newValue, oldValue, url }) => {
  if (key === "carnet:articles") rendre(JSON.parse(newValue ?? "[]"));
});
\end{codejs}

\section{Cookies~: le seul stockage vu par le serveur}

\begin{codejs}
// Ecriture depuis JS : pas de HttpOnly (sinon JS ne peut pas le lire), 4 Ko max
document.cookie = "theme=sombre; path=/; max-age=31536000; SameSite=Lax; Secure";

const cookies = Object.fromEntries(
  document.cookie.split("; ").map((p) => { const [k, ...v] = p.split("="); return [k, decodeURIComponent(v.join("="))]; })
);
const effacer = (nom) => { document.cookie = `${nom}=; path=/; max-age=0`; };
\end{codejs}

\begin{tabz}{lX}
  \headrow \thd{Attribut} \& \thd{Effet} \\
  \midrule
  \verb|Expires| / \verb|Max-Age| \& dur\'ee de vie~; sans lui, cookie de session (mort \`a la fermeture). \\
  \verb|Path| \& ne sera envoy\'e que sur ce chemin. \verb|/| dans 99 \% des cas. \\
  \verb|Domain| \& sous-domaines inclus ; ne le mets pas sauf besoin r\'eel (\verb|.exemple.com| = large surface). \\
  \verb|Secure| \& uniquement en HTTPS --- obligatoire si \verb|SameSite=None|. \\
  \verb|HttpOnly| \& \textbfJavaScript ne peut pas le lire~: la meilleure d\'efense contre le vol de session par XSS. \\
  \verb|SameSite=Lax|  \& d\'efaut moderne~; bloque la plupart des CSRF. \verb|Strict| = encore plus raide. \verb|None| = tiers, exige \verb|Secure|. \\
  \verb|Partitioned| \& cookie cloisonn\'e par site parent (CHIPS) : pour les iframe tierces. \\
\end{tabz}

\begin{note}
Depuis quelques ann\'ees, les navigateurs traitent un cookie \verb|SameSite=None| sans
\verb|Secure| comme invalide, et bloquent les cookies tiers par d\'efaut. Conclusion~: une
session doit \^etre un cookie \verb|HttpOnly; Secure; SameSite=Lax|, et ton JavaScript ne doit
jamais avoir besoin de lire le jeton de session. Si tu stockes un jeton d'acc\`es dans
\verb|localStorage| pour \'eviter cela, accepte d'\^etre vuln\'erable au XSS~: c'est le vrai d\'ebat
JWT-localStorage vs cookie-HttpOnly, et la r\'eponse raisonnable est \og ~cookie + CSRF~ \fg{}.
\end{note}

\section{IndexedDB --- la base de donn\'ees du navigateur}

C'est une base orient\'ee enregistrements~: magasins (\texttt{objectStore}), cl\'es, index,
transactions. API de bas niveau verbeuse~; \link{idb}{https://github.com/jakearchibald/idb}
(\cmd{npm i idb}) la rend supportable sans lui retirer sa puissance.

\begin{codejs}
import { openDB } from "idb";

const db = await openDB("carnet", 2, {
  upgrade(bdd, ancienne, nouvelle, t) {
    if (!bdd.objectStoreNames.contains("articles")) {
      const s = bdd.createObjectStore("articles", { keyPath: "id", autoIncrement: true });
      s.createIndex("par-date", "creeLe");
      s.createIndex("par-auteur-\'etat", ["auteur", "etat"], { multiEntry: false });
      s.createIndex("pleine-texte", "recherche", { multiEntry: true });   // tableau de mots-cles
    }
    if (nouvelle === 2 && !bdd.objectStoreNames.contains("tags")) {
      bdd.createObjectStore("tags", { keyPath: "nom" });
    }
  },
});

// Ecrire
const id = await db.put("articles", { titre: "Apprendre JS", etat: "brouillon", creeLe: Date.now(), recherche: ["apprendre", "js"] });
await db.add("articles", { titre: "unique" });           // echoue si la cle existe
await db.bulkPut("articles", tableauDObjects);
await db.delete("articles", id);
await db.clear("articles");

// Lire
const un = await db.get("articles", id);
const tous = await db.getAll("articles");
const compte = await db.count("articles");
const parDate = await db.getAllFromIndex("articles", "par-date", IDBKeyRange.bound(d1, d2));
const cursor = await db.transaction("articles").store.index("par-date").openCursor(IDBKeyRange.lowerBound(Date.now() - 864e5), "prev");
const page = [];
while (cursor && page.length < 20) { page.push(cursor.value); cursor = await cursor.continue(); }
const cleSuivante = await db.getAllKeys("articles", null, 1);   // pagination par cle, la scalable

// Transaction multi-magasins, atomique
const t = db.transaction(["articles", "tags"], "readwrite");
await Promise.all([
  t.objectStore("articles").put(article),
  t.objectStore("tags").put({ nom: "js", nb: (await t.objectStore("tags").get("js"))?.nb + 1 ?? 1 }),
]);
await t.done;      // commit ; "aborted" en cas d'erreur -> tout est annule

// Recherche pleine texte maison (index multiEntry) :
const mots = requete.toLowerCase().split(/\s+/).filter(Boolean);
const resultat = await db.getAllFromIndex("articles", "pleine-texte", IDBKeyRange.only(mots[0]));
\end{codejs}

\begin{astuce}
Besoin de requ\^etes SQL, de jointures et de gros volumes dans le navigateur~?
\link{SQL.js}{https://sql.js.org/}, \link{wa-sqlite}{https://github.com/rhashimoto/wa-sqlite} ou
\link{TinyBase}{https://tinybase.org/} tournent sur IndexedDB/OPFS. Pour une application locale
lourde (\'editeur, cartographie, comptabilit\'e), SQLite compil\'e en WASM sur OPFS est
une architecture r\'eelle, pas un gadget.
\end{astuce}

\section{Service worker, cache et hors ligne}
\label{sec:pwa}

Un \emph{service worker} est un script qui s'installe dans le navigateur, hors de la page,
et qui peut intercepter les requ\^etes r\'eseau. C'est ce qui rend un site installable, rapide en
2\textsuperscript{e} visite, et utilisable sans r\'eseau.

\begin{codejs}
// sw.js --- a la racine du site (portee = son dossier et au-dessous)
const CACHE = "v3";
const ESSENTIELS = ["/", "/index.html", "/app.js", "/style.css", "/icone-192.png"];

self.addEventListener("install", (e) => {
  e.waitUntil(caches.open(CACHE).then((c) => c.addAll(ESSENTIELS)).then(() => self.skipWaiting()));
});
self.addEventListener("activate", (e) => {
  e.waitUntil(
    caches.keys()
      .then((cles) => Promise.all(cles.filter((k) => k !== CACHE).map((k) => caches.delete(k))))
      .then(() => self.clients.claim())
  );
});
self.addEventListener("fetch", (e) => {
  const req = e.request;
  if (req.method !== "GET" || new URL(req.url).origin !== self.location.origin) return;

  // 1) Resources applicatives : stale-while-revalidate
  if (req.destination === "script" || req.destination === "style" || req.mode === "navigate") {
    e.respondWith(
      caches.match(req).then((cache) => {
        const reseau = fetch(req)
          .then((res) => { const clone = res.clone(); caches.open(CACHE).then((c) => c.put(req, clone)); return res; })
          .catch(() => cache ?? Response.error());
        return cache ?? reseau;
      })
    );
    return;
  }
  // 2) API : network-first avec repli sur le cache (donnees legerement perimees > rien)
  if (new URL(req.url).pathname.startsWith("/api/")) {
    e.respondWith(
      fetch(req)
        .then((res) => { if (res.ok) caches.open(CACHE).then((c) => c.put(req, res.clone())); return res; })
        .catch(() => caches.match(req).then((r) => r ?? new Response(JSON.stringify({ horsLigne: true }), { headers: { "content-type": "application/json" } })))
    );
  }
});
// Recevoir un message de la page pour demander une mise a jour
self.addEventListener("message", (e) => { if (e.data?.type === "verifier") self.skipWaiting(); });
\end{codejs}

\begin{codejs}
// enregistrement depuis la page
if ("serviceWorker" in navigator && location.protocol === "https:") {
  const reg = await navigator.serviceWorker.register("/sw.js", { scope: "/" });
  reg.addEventListener("updatefound", () => {
    const nouveau = reg.installing;
    nouveau?.addEventListener("statechange", () => {
      if (nouveau.state === "installed" && navigator.serviceWorker.controller) {
        notifier("Nouvelle version disponible", { action: { titre: "Recharger", clic: () => nouveau.postMessage({ type: "aktiver" }) } });
      }
    });
  });
}
// Detection d'etat reseau
const enLigne = navigator.onLine;
window.addEventListener("online", reSynchroniser);
window.addEventListener("offline", () => afficherBanniere("Mode hors ligne : tes saisies seront renvoyees plus tard"));
// File d'attente d'arriere-plan officielle (Chromium + Firefox 2024+)
if ("BackgroundSync" in window) { const reg = await navigator.serviceWorker.getRegistration(); await reg.sync.register("renvoi"); }
\end{codejs}

\begin{tabz}{lX}
  \headrow \thd{Strat\'egie de cache} \& \thd{Quand} \\
  \midrule
  Cache-only \& ressources fig\'ees, versionn\'ees dans le nom. \\
  Network-only \& appels non-idempotents, endpoints sensibles. \\
  Network-first, cache en repli \& donn\'ees qui doivent \^etre fra\^iches (API, HTML). \\
  Cache-first, r\'eseau en repli \& assets immuables, images. \\
  Stale-while-revalidate \& le meilleur compromis pour une SPA~: affichage imm\'ediat + mise \`a jour en fond. \\
  Race \& deux miroirs, on garde le premier arriv\'e (langueur r\'eseau \'elev\'ee). \\
\end{tabz}

\etape{Rendre la PWA installable}
\begin{codejson}
{
  "name": "Carnet de Nour",
  "short_name": "Carnet",
  "id": "/?source=pwa",
  "start_url": "/",
  "scope": "/",
  "display": "standalone",
  "display_override": ["window-controls-overlay"],
  "orientation": "any",
  "theme_color": "#0b4f8a",
  "background_color": "#ffffff",
  "description": "Prise de notes rapide, hors ligne, chiffre cote appareil.",
  "icons": [
    { "src": "/icones/icone-192.png", "sizes": "192x192", "type": "image/png", "purpose": "any maskable" },
    { "src": "/icones/icone-512.png", "sizes": "512x512", "type": "image/png" }
  ],
  "shortcuts": [{ "name": "Nouvelle note", "url": "/?nouvelle", "icons": [{ "src": "/icones/note.png", "sizes": "96x96" }] }]
}
\end{codejson}

\begin{terminal}
npm i -D vite-plugin-pwa       # genere manifest, sw, icone, et le cycle de mise a jour
npx @vite/pwa-assets-creator   # genere toutes les tailles d'icone depuis un SVG
npx lighthouse --only-categories=pwa --output=html --output-path=rapport.html http://localhost:4173
\end{terminal}

\section{S\'ecurit\'e et confidentialit\'e du stockage}

\begin{itemize}
  \item Ne chiffre rien \og ~\`a la main~ \fg{} dans le code client en pensant que c'est secret~:
        la clef est dans le bundle. Pour du secret \emph{par appareil}, la WebCrypto API avec une
        clef non extractable (\verb|CryptoKey| \verb|extractable: false|) est correcte~; le reste
        est une d\'eclaration d'intention.
  \item Limite la r\'etention~: purge les cl\'es plus vieilles que X jours, vide le stockage \`a la
        d\'econnexion (\verb|IndexedDB.deleteDatabase|, purge du cache, nettoyage SW).
  \item RGPD~: un stockage local non-transmis reste une donn\'ee personnelle~; documente-le dans la
        politique de confidentialit\'e.
  \item \link{OPFS}{https://developer.mozilla.org/en-US/docs/Web/API/File\_System\_API/Origin\_private\_file\_system}
        est priv\'e \`a l'origine~: plus rapide et plus gros que IndexedDB pour des fichiers.
\end{itemize}

\begin{terminal}
# vider tout ce que le site a garde (DevTools > Application > Storage)
const cachesKeys = await caches.keys(); await Promise.all(cachesKeys.map(k => caches.delete(k)));
indexedDB.databases().then(dbs => dbs.forEach(d => indexedDB.deleteDatabase(d.name)));
localStorage.clear(); sessionStorage.clear();
navigator.serviceWorker.getRegistrations().then(rs => rs.forEach(r => r.unregister()));
location.reload();
\end{terminal}

\begin{exercice}
\begin{enumerate}
  \item Ajoute \`a ta liste de courses une sauvegarde automatique~: apr\`es chaque mutation,
        \'ecrit dans \verb|localStorage| avec un \verb|debounce| de 400~ms, et restaure au
        d\'emarrage en affichant \og ~brouillon restaur\'e~ \fg{}.
  \item \'Ecris une fonction de \textbf{synchronisation}~: un tableau \verb|enAttente| dans
        IndexedDB, renvoi automatique quand \verb|online| revient, avec gestion des doublons
        par identifiant local.
  \item Migre ton stockage de \verb|localStorage| vers IndexedDB~: garde une cl\'e de version,
        et un code qui lit l'ancien format une seule fois puis le supprime.
  \item Ajoute une exportation JSON (t\'el\'echargement du fichier via Blob) et une importation
        valid\'ee par \verb|zod|~: refuse tout fichier malform\'e sans corrompre l'\'etat existant.
    \item Rends ton application installable~: manifest, service worker, mode hors ligne
        complet, et un indicateur visible \og ~derni\`ere synchronisation~ \fg{}.
  \item V\'erifie dans DevTools \$\rightarrow\$ Application \$\rightarrow\$ Storage ce que ton site
        garde r\'eellement~; supprime ce qui ne devrait pas y \^etre.
\end{enumerate}
\end{exercice}

\begin{ressources}
\begin{itemize}
  \item MDN --- Web Storage --- \url{https://developer.mozilla.org/fr/docs/Web/API/Window/localStorage}
  \item MDN --- IndexedDB --- \url{https://developer.mozilla.org/fr/docs/Web/API/IndexedDB\\\_API}
  \item \emph{The Offline Cookbook} (Jake Archibald) --- \url{https://web.dev/articles/offline-cookbook}
  \item idb --- \url{https://github.com/jakearchibald/idb}
  \item PWA --- \url{https://web.dev/learn/pwa/} --- Workbox --- \url{https://developer.chrome.com/docs/workbox}
\end{itemize}
\end{ressources}

\index{IndexedDB}\index{service worker}\index{PWA}\index{cookies}
